\documentclass[11pt]{article}
\usepackage[T1]{fontenc}
\usepackage[utf8]{inputenc}
\usepackage{lmodern,amsmath,amssymb,amsthm,microtype,booktabs}
\usepackage[a4paper,margin=28mm]{geometry}
\usepackage[colorlinks=true,linkcolor=blue,citecolor=blue,urlcolor=blue]{hyperref}
\hypersetup{
  pdftitle={Growth criteria for reciprocal series with nonincreasing weights},
  pdfauthor={twinheterodox},
  pdfsubject={Irrationality of weighted reciprocal series under irregular growth},
  pdfcreator={LaTeX}
}
\newtheorem{theorem}{Theorem}[section]
\newtheorem{lemma}[theorem]{Lemma}
\newtheorem{proposition}[theorem]{Proposition}
\numberwithin{equation}{section}
\newcommand{\Q}{\mathbb Q}
\newcommand{\N}{\mathbb N}
\newcommand{\bw}{\mathbf w}
\title{Growth criteria for reciprocal series with nonincreasing weights}
\author{twinheterodox\\[2pt]\normalsize Prepared with Codex}
\date{13 November 2026}

\begin{document}
\maketitle

\begin{abstract}
Fix nonincreasing positive integer weights $w_0,\ldots,w_{d-1}$.
Define $c>1$ by
\[
  (c-1)^2\sum_{j=0}^{d-1}w_jc^j=w_0.
\]
If the total weight is at least two, we prove irrationality of
$\sum_{n\ge1}1/(a_n^{w_0}\cdots a_{n+d-1}^{w_{d-1}})$ whenever the
positive integers $a_n$ are strictly increasing and
\mbox{$\limsup_{n\to\infty}(\log a_n)/c^n=\infty$}.
For every nonconstant nonincreasing weight profile, this lowers the
sufficient base in the general criterion of Barreto, Kang, Kim,
Kova\v{c} and Zhang \cite[Theorem~3]{BKKKZ}, with the classical
base-$2$ criterion included in the comparison.
For weights $(3,2,1)$ the base becomes $\sqrt[3]{4}$;
staircase weights $(d,d-1,\ldots,1)$ give $(d+1)^{1/d}$.
Irregular growth is allowed. A more general form permits nondecreasing
sequences and positive integer numerators under polynomial size bounds.
\end{abstract}

\section{Introduction}

Growth criteria for irrationality apply to families of reciprocal sums
without requiring a closed form for the sum or information about the
prime factors of the denominators. Their arithmetic starting point is
simple: if the whole sum were rational, a positive tail could not become
arbitrarily small after clearing the preceding partial sum. For products
of overlapping sequence terms, the powers needed to clear that partial
sum depend on the position of each factor. Keeping this dependence
improves the growth condition.

Barreto, Kang, Kim, Kova\v{c} and Zhang \cite{BKKKZ} answered a question
of Erd\H{o}s and Graham \cite[p.~64]{ErdosGraham} about reciprocal
products of consecutive terms and established criteria for general
weight profiles. Their sufficient conditions and rational constructions
leave a gap for many profiles with repeated factors
\cite[Section~5]{BKKKZ}. An earlier paper \cite{Earlier} improved their
sufficient base from $2$ to $\sqrt3$ for weights $(2,1)$.
Here we extend the endpoint argument to every nonincreasing positive
weight profile. A positive combination of cumulative logarithms is
what makes the extension work for irregular growth.

\subsection{The result}

Throughout the paper, $d\ge1$ and
\begin{equation}\label{eq:weights}
  w_0\ge w_1\ge\cdots\ge w_{d-1}\ge1
\end{equation}
are fixed integers. Write $s=\sum_{j=0}^{d-1}w_j$ and
$f_{\bw}(x)=\sum_{j=0}^{d-1}w_jx^j$. Define $c=c_{\bw}>1$ by
\begin{equation}\label{eq:root}
  (c-1)^2 f_{\bw}(c)=w_0.
\end{equation}
The left-hand side is strictly increasing from zero to infinity for
$c>1$, so the definition is unambiguous. All logarithms are natural.

\begin{theorem}\label{thm:strict}
Suppose $s\ge2$. If $(a_n)_{n\ge1}$ is a strictly increasing sequence
of positive integers and
\begin{equation}\label{eq:growth}
  \limsup_{n\to\infty}\frac{\log a_n}{c^n}=\infty,
\end{equation}
then
\begin{equation}\label{eq:unitseries}
  \sum_{n=1}^{\infty}
  \frac{1}{a_n^{w_0}a_{n+1}^{w_1}\cdots a_{n+d-1}^{w_{d-1}}}
\end{equation}
is irrational.
\end{theorem}

Condition~\eqref{eq:growth} is equivalent to
$\limsup a_n^{1/c^n}=\infty$. It permits slow stretches and abrupt
jumps: no ratio limit for successive logarithms is assumed.
We prove the following form, which also permits repetitions and
positive integer numerators.

\begin{theorem}\label{thm:general}
Let $(a_n)$ be nondecreasing positive integers and $(b_n)$ positive
integers. Suppose that, for some real numbers $0\le\eta<\tau$,
\begin{equation}\label{eq:size}
  b_n\le n^\eta,\qquad
  x_n:=\prod_{j=0}^{d-1}a_{n+j}^{w_j}\ge n^{1+\tau}
\end{equation}
for all sufficiently large $n$. If \eqref{eq:growth} holds, then the
convergent series $\sum_{n\ge1}b_n/x_n$ is irrational.
\end{theorem}

For Theorem~\ref{thm:strict}, take $b_n=1$, $\eta=0$ and $\tau=s-1$,
using $a_n\ge n$. The case $d=1,w_0=1$ has $c=2$ and retains the
size hypothesis in Theorem~\ref{thm:general}; strict increase alone is
not substituted for that hypothesis. The polynomial bounds in
\eqref{eq:size} control tails where the sequence grows slowly.

\subsection{What improves}

Let $\rho_{\bw}$ be the positive root used in
\cite[Theorem~3]{BKKKZ}. For the profiles \eqref{eq:weights}, it solves
\begin{equation}\label{eq:oldroot}
  (x-1)f_{\bw}(x)-w_0x^{d-1}=0.
\end{equation}
The classical base-$2$ criterion of Erd\H{o}s \cite[Theorem~1]{Erdos1975}
also applies to the integer product denominators with unit numerators;
the same fallback for the general setting follows from
\cite[Theorem~3, $d=1$]{BKKKZ}. Thus the relevant comparison is with
$\min(\rho_{\bw},2)$, as emphasized in \cite[Remark~4(4)]{BKKKZ}.
We prove in Proposition~\ref{prop:roots} that
\begin{equation}\label{eq:improvement}
  c_{\bw}<\min(\rho_{\bw},2)
\end{equation}
for every nonconstant profile satisfying \eqref{eq:weights}.
Constant profiles recover the earlier base.

\begin{center}
\renewcommand{\arraystretch}{1.35}
\begin{tabular}{lcc}
\toprule
Weights & Previous base $\min(\rho_{\bw},2)$ & Base in this paper\\
\midrule
$(2,1)$ & $2$ & $\sqrt3$\\
$(3,2)$ & $(1+\sqrt7)/2$ & $(1+\sqrt{33})/4$\\
$(3,2,1)$ & $2$ & $\sqrt[3]{4}$\\
\bottomrule
\end{tabular}
\end{center}

The first row recovers \cite{Earlier}, as do the profiles $(2k,k)$
after replacing $a_n$ by $a_n^k$. For staircase weights
$(d,d-1,\ldots,1)$, equation~\eqref{eq:root} reduces to
\begin{equation}\label{eq:staircase}
  c=(d+1)^{1/d}.
\end{equation}
For example, $a_n=\lceil\exp((8/5)^n)\rceil$ is covered for weights
$(3,2,1)$ because $8/5>\sqrt[3]{4}$, while $(\log a_n)/2^n\to0$.
More generally, any regular growth rate $\log a_n\sim\theta^n$ with
$c_{\bw}<\theta<\min(\rho_{\bw},2)$ distinguishes the two criteria.
These examples illustrate the improvement; the theorem does not
require regular growth. The optimal thresholds remain open in general.

\subsection{Why the endpoint matters}

For weights $(3,2,1)$, the first $N$ terms are cleared by
\[
  D_N=\left(\prod_{i=1}^{N}a_i^3\right)a_{N+1}^2a_{N+2}.
\]
The last two factors need powers two and one, respectively, rather
than three. With $u_n=\log a_n$ and $V_N=\sum_{i=1}^N u_i$, this gives
\[
  \log\frac{D_N}{a_{N+1}^3a_{N+2}^2a_{N+3}}
  =4V_N-V_{N+3}.
\]
The recurrence $V_{N+3}=4V_N$ explains the cube-root base. The harder
question is whether the saving persists for the entire tail when
growth is irregular. For general nonincreasing weights, we construct
a positive combination $W_N$ of $V_N,\ldots,V_{N+d-1}$ for which the
same comparison becomes $cW_N-W_{N+1}$. A carefully selected first
peak then controls the near tail, and a later large term controls the
far tail using the same denominator.

We adapt the positive-integer tail argument, uniform tail estimates,
peak selection and fast/slow splitting in
\cite[Lemmas~8--10 and Proposition~12]{BKKKZ}, whose antecedents include
Erd\H{o}s's arguments \cite[pp.~3--6]{Erdos1975}.
The endpoint refinement and cumulative-log approach for $(2,1)$
come from \cite{Earlier}. The new ingredient is the positive
factorization for every profile \eqref{eq:weights}, together with
control of its forward shifts at the last slow index.
The proof below supplies all required estimates.

\section{Uniform tail estimates}\label{sec:tails}

We need one estimate valid under polynomial growth and a sharper one
on intervals of exponential growth. Both intervals will move during
the proof, so their constants must be independent of the endpoints.
These are the relevant forms of \cite[Lemma~10]{BKKKZ}.

\begin{lemma}\label{lem:tails}
Let $(x_n)$ be nondecreasing positive real numbers, and fix
$0\le\eta<\tau$. Let $A$ be a positive integer and suppose
$0\le y_n\le n^\eta$ on the indices being summed.
\begin{enumerate}
\item If $x_n\ge n^{1+\tau}$ for all $n\ge A$, then
\begin{equation}\label{eq:polytail}
  \sum_{n\ge A}\frac{y_n}{x_n}
  \le C_{\mathrm{pol}}x_A^{-\alpha},
  \qquad \alpha=\frac{\tau-\eta}{1+\tau}>0.
\end{equation}
Here $C_{\mathrm{pol}}$ depends only on $\eta,\tau$.
\item If $x_n\ge e^n$ for $A\le n\le B$, where $B\ge A$ is an integer
or $B=\infty$, then
\begin{equation}\label{eq:exptail}
  \sum_{n=A}^{B}\frac{y_n}{x_n}
  \le C_{\mathrm{exp}}\frac{(\log x_A)^{1+\eta}}{x_A}.
\end{equation}
Here $C_{\mathrm{exp}}$ depends only on $\eta$.
\end{enumerate}
Neither constant depends on the endpoints or on the sequences. No
numerator bound before $A$ is required.
\end{lemma}

\begin{proof}
Put $X=x_A$ and group indices by $2^jX\le x_n<2^{j+1}X$, $j\ge0$.
In the first case every index in the $j$th group is at most
$(2^{j+1}X)^{1/(1+\tau)}$. Summing $n^\eta$ up to that bound shows
that the group's contribution is at most
\[
  C\,\frac{(2^{j+1}X)^{(1+\eta)/(1+\tau)}}{2^jX}
  \le C'X^{-\alpha}2^{-\alpha j}.
\]
The geometric series converges, proving \eqref{eq:polytail}.
In the second case, every index in the $j$th group is at most
$\log X+(j+1)\log2$. Its contribution is at most
\[
  C\,\frac{(\log X+j+1)^{1+\eta}}{2^jX}.
\]
Since $\log X\ge A\ge1$, summing over $j$ proves
\eqref{eq:exptail}. The same estimate applies if a finite interval
ends partway through a group.
\end{proof}

\section{Proof of the irrationality criterion}\label{sec:proof}

\subsection{Endpoint clearing and a positive combination}

Retain $x_n$ from \eqref{eq:size}, and put $u_n=\log a_n$, $V_0=0$,
$V_N=\sum_{i=1}^N u_i$. For $N\ge1$, define
\begin{equation}\label{eq:clearer}
  D_N=\left(\prod_{i=1}^{N}a_i^{w_0}\right)
          \prod_{r=1}^{d-1}a_{N+r}^{w_r},\qquad
  E_N=\log D_N,\qquad F_N=\log x_{N+1}.
\end{equation}
The integer $D_N$ clears every term through $N$. Factors of index at
most $N$ are allotted the largest weight $w_0$. At index $N+r>N$,
an earlier term can use only a weight $w_j$ with $j\ge r$, hence
$w_j\le w_r$. Thus divisibility holds even when some $a_i$ coincide
or share prime factors. We do not claim that $D_N$ is reduced.

Set $w_d=0$ and
\[
  \lambda_r=w_{r-1}-w_r\quad(1\le r\le d).
\]
Then $\lambda_r\ge0$, $\lambda_d=w_{d-1}>0$, and direct subtraction
gives
\begin{align}
  E_N-F_N&=w_0V_N-\sum_{r=1}^d\lambda_r u_{N+r},
       \label{eq:endpoint}\\
  E_{N+1}-E_N&=\sum_{r=1}^d\lambda_r u_{N+r}\le F_N.
       \label{eq:increment}
\end{align}
To use \eqref{eq:endpoint} at peaks, we need a positive combination
of cumulative logarithms. Introduce
\[
  q_{\bw}(x)=(x-1)\sum_{r=1}^d\lambda_r x^{r-1}-w_0.
\]
The identity
\begin{equation}\label{eq:rootidentity}
  xq_{\bw}(x)=(x-1)^2f_{\bw}(x)-w_0
\end{equation}
shows that $q_{\bw}(c)=0$ for the base in \eqref{eq:root}.
Division by this root has a useful sign property:
\begin{equation}\label{eq:filter}
\begin{split}
  q_{\bw}(x)&=(x-c)H(x),\\
  H(x)&=\sum_{r=1}^d\lambda_r\left[
    x^{r-1}+(c-1)\sum_{k=0}^{r-2}c^{r-2-k}x^k\right]
       =\sum_{j=0}^{d-1}h_jx^j.
\end{split}
\end{equation}
The inner sum is empty for $r=1$. To verify the formula, subtract
the value at $c$ from each $(x-1)x^{r-1}$, divide by $x-c$, and use
$q_{\bw}(c)=0$. Every $h_j$ is positive: the term with $r=d$ supplies
all lower degrees, and $h_{d-1}=\lambda_d=w_{d-1}>0$.
For $d=1$ the formula simply gives $H=w_0$.

Define
\begin{equation}\label{eq:WZ}
  W_N=\sum_{j=0}^{d-1}h_jV_{N+j},\qquad Z_N=\frac{W_N}{c^N}.
\end{equation}
Substituting $u_{N+r}=V_{N+r}-V_{N+r-1}$ in
\eqref{eq:endpoint} and using \eqref{eq:filter} gives the exact identity
\begin{equation}\label{eq:positive}
  E_N-F_N=cW_N-W_{N+1}.
\end{equation}
This identity preserves the endpoint saving and allows us to compare
positive quantities at successive peaks.

We also need to control the forward shifts in $W_N$. Since all $u_i$
are nonnegative and $V_i$ is nondecreasing, with $H_*:=\sum_jh_j$,
\begin{align}
  w_{d-1}V_{N+d-1}&\le E_N\le w_0V_{N+d-1},\notag\\
  w_{d-1}V_{N+d-1}&\le W_N\le H_*V_{N+d-1}. \label{eq:compare}
\end{align}
For the first line one may also use
$E_N=\sum_r\lambda_rV_{N+r-1}$.
Writing $C_E=w_0/w_{d-1}$ and $C_W=H_*/w_{d-1}$, we obtain
\begin{equation}\label{eq:EW}
  E_N\le C_EW_N,\qquad W_N\le C_WE_N,\qquad
  \log x_N\le w_0V_{N+d-1}\le C_EW_N.
\end{equation}
These constants depend only on the fixed weights.

\subsection{The jump bound forced by rationality}

Suppose for a contradiction that $S=\sum b_n/x_n=p/q$, where $q$ is a
positive integer. The series converges by \eqref{eq:size}. Put
$T_N=\sum_{n>N}b_n/x_n$. Since $D_N$ clears the first $N$ terms and
the tail is positive, $qD_NT_N$ is a positive integer. Thus
\begin{equation}\label{eq:integerlower}
  D_NT_N\ge 1/q \qquad(N\ge1).
\end{equation}
This is the elementary criterion in \cite[Lemma~8]{BKKKZ}.

The denominators $x_n$ are nondecreasing. Applying
\eqref{eq:polytail} to $T_N$ and combining it with
\eqref{eq:integerlower} yields, for all sufficiently large $N$,
\[
  \alpha F_N\le E_N+\log(C_{\mathrm{pol}}q).
\]
Monotonicity and \eqref{eq:size} imply $a_n\to\infty$, so $E_N\to\infty$.
Consequently,
\begin{equation}\label{eq:jump}
  F_N\le \frac2\alpha E_N,\qquad
  E_{N+1}\le A E_N,\qquad A=1+\frac2\alpha
\end{equation}
for all large $N$, the second inequality following from
\eqref{eq:increment}. This coarse bound prevents an uncontrolled jump
immediately after a slow index under the rationality assumption.
All constants below may depend on the fixed weights, $\eta,\tau$,
the sequence and $q$, but not on the indices chosen later.

\subsection{Peaks and eventual exponential growth}

By \eqref{eq:growth} and $h_0>0$, the sequence $Z_N$ is unbounded.
There are therefore infinitely many indices $Q$ such that
\begin{equation}\label{eq:globalpeak}
  Z_Q>(1+Q^{-2})\max(1,Z_1,\ldots,Z_{Q-1}).
\end{equation}
Otherwise the successive running maxima would eventually be bounded
by a fixed constant times
\begin{equation}\label{eq:peakproduct}
  \prod_{j\ge1}(1+j^{-2})
  \le\exp\left(\sum_{j\ge1}j^{-2}\right)\le e^2.
\end{equation}
This is the summable-increment peak argument used in
\cite[Lemma~9]{BKKKZ}. At any such peak, \eqref{eq:positive} gives
\begin{equation}\label{eq:saving}
  E_{Q-1}-F_{Q-1}=cW_{Q-1}-W_Q
       \le-\frac{W_Q}{Q^2+1},\qquad W_Q>c^Q.
\end{equation}

Suppose first that $a_n\ge e^n$ for every sufficiently large $n$.
Then $x_n\ge a_n^{w_0}\ge e^n$ there. At large peaks $Q$,
\eqref{eq:exptail}, \eqref{eq:saving} and \eqref{eq:EW} give
\[
  \log(D_{Q-1}T_{Q-1})
  \le-\frac{W_Q}{Q^2+1}+(1+\eta)\log(C_EW_Q)+O(1).
\]
For large $Q$, the right-hand expression is a decreasing function of
$W_Q$ throughout $W_Q\ge c^Q$. It is therefore at most
\[
  -\frac{c^Q}{Q^2+1}+(1+\eta)(Q\log c+\log C_E)+O(1),
\]
which tends to minus infinity. This contradicts
\eqref{eq:integerlower}. Notice that no upper bound on the size of a
peak was needed in this case.

\subsection{A first peak after a slow index}

It remains to allow $a_n<e^n$ at arbitrarily large indices. Choose $R$
tending to infinity with
\begin{equation}\label{eq:mu}
  \mu_R=\frac{u_R}{c^R}\longrightarrow\infty.
\end{equation}
For large $R$ one has $a_R\ge e^R$. Let $P=P(R)<R$ be the last index
before $R$ with $a_P<e^P$. For every fixed slow index $t$, the selected
$R$ eventually exceed $t$, and then $P(R)\ge t$. Thus $P(R)\to\infty$,
and
\begin{equation}\label{eq:faststretch}
  a_n\ge e^n\qquad(P<n\le R).
\end{equation}

The quantity $W_P$ involves cumulative sums as far as $P+d-1$, whose
size is not supplied by $a_P<e^P$. This is where the jump estimate
\eqref{eq:jump} is needed. All factors of $D_{P-d+1}$ have index at
most $P$, so, once $P$ is large enough,
\[
  E_{P-d+1}\le w_0V_P\le w_0P u_P<w_0P^2.
\]
Applying \eqref{eq:jump} exactly $d-1$ times, and then \eqref{eq:EW}, gives
\begin{equation}\label{eq:slowcontrol}
  E_P\le w_0A^{d-1}P^2,\qquad
  W_P\le C_Ww_0A^{d-1}P^2.
\end{equation}
For $d=1$ there are no jump steps. In every case $Z_P<1$ for large $P$.

Put $J=R-d+1$. The last term in $W_J$ gives
\begin{equation}\label{eq:largeendpoint}
  Z_J\ge\frac{w_{d-1}u_R}{c^{R-d+1}}
       =w_{d-1}c^{d-1}\mu_R\longrightarrow\infty.
\end{equation}
We have $J>P$ eventually. Indeed, otherwise $P<R\le P+d-1$, so
$W_P\ge w_{d-1}u_R$. Equation~\eqref{eq:slowcontrol} would imply
$\mu_R\le CP^2/c^R\to0$, contradicting \eqref{eq:mu}.

Choose the first $Q\in[P+1,J]$ for which
\begin{equation}\label{eq:firstpeak}
  Z_Q>(1+Q^{-2})\max(1,Z_P,\ldots,Z_{Q-1}).
\end{equation}
It exists for all large $R$: if no such peak occurred, the product
bound \eqref{eq:peakproduct} would give $Z_J\le e^2$, contrary to
\eqref{eq:largeendpoint}. Because this is the first peak, all preceding
running maxima are bounded by $L=e^2$. In particular,
\begin{equation}\label{eq:beforepeak}
  E_{Q-1}\le C_E L c^Q,\qquad
  W_Q>c^Q,\qquad Q\le R-d+1\le R.
\end{equation}
Also $Q\to\infty$, since $Q\ge P+1$ and $P\to\infty$.
The saving \eqref{eq:saving} remains valid because $Q-1\ge P$.

\subsection{Controlling the same cleared tail on both sides}

Split the tail at $R$, retaining one multiplier:
\begin{equation}\label{eq:split}
  D_{Q-1}T_{Q-1}
   =D_{Q-1}\sum_{n=Q}^{R}\frac{b_n}{x_n}+D_{Q-1}T_R.
\end{equation}
For the near part, $x_n\ge a_n^{w_0}\ge e^n$ on $Q\le n\le R$.
Only this first factor is needed; other factors of $x_n$ may lie
beyond the fast stretch. By \eqref{eq:jump} and \eqref{eq:beforepeak},
\[
  F_{Q-1}\le\frac{2C_E L}{\alpha}c^Q.
\]
The finite-interval estimate \eqref{eq:exptail} and the saving
\eqref{eq:saving} now imply
\begin{align}
  \log\left(D_{Q-1}\sum_{n=Q}^{R}\frac{b_n}{x_n}\right)
  &\le E_{Q-1}-F_{Q-1}+(1+\eta)\log F_{Q-1}+O(1)\notag\\
  &\le-\frac{c^Q}{Q^2+1}+(1+\eta)Q\log c+O(1)
       \longrightarrow-\infty. \label{eq:near}
\end{align}
The constants are independent of $P,Q,R$.

For the far part, monotonicity gives $x_{R+1}\ge a_R^s$.
Using \eqref{eq:polytail}, \eqref{eq:beforepeak} and \eqref{eq:mu},
\begin{align}
  D_{Q-1}T_R
   &\le C_{\mathrm{pol}}\exp(E_{Q-1}-\alpha\log x_{R+1})\notag\\
   &\le C_{\mathrm{pol}}\exp(C_E Lc^Q-\alpha s u_R)\notag\\
   &\le C_{\mathrm{pol}}\exp\bigl((C_E L-\alpha s\mu_R)c^R\bigr)
        \longrightarrow0. \label{eq:far}
\end{align}
Both terms in \eqref{eq:split} tend to zero along the same selected
indices, contradicting \eqref{eq:integerlower}. This proves
Theorem~\ref{thm:general} and hence Theorem~\ref{thm:strict}. \qed

\section{Comparison of the growth bases}\label{sec:comparison}

The comparisons below separate the improvement in the sufficient base
from the question of sharpness. They also show that a natural class of
product substitutions into the earlier theorem does not recover the
new base.

\begin{proposition}\label{prop:roots}
Let $\psi_d>1$ solve $\psi_d^d-\psi_d^{d-1}=1$, and let
$\rho_{\bw}$ be the root in \eqref{eq:oldroot}.
For constant profiles, $c_{\bw}=\rho_{\bw}=\psi_d$.
For nonconstant profiles satisfying \eqref{eq:weights},
\[
  \psi_d<c_{\bw}<\min(\rho_{\bw},2).
\]
\end{proposition}

\begin{proof}
Write $T_d(x)=1+x+\cdots+x^{d-1}$. For $x>1$,
$f_{\bw}(x)/w_0\le T_d(x)$, strictly unless the profile is constant.
At $x=\psi_d$, one has $(x-1)^2T_d(x)=1$.
The monotonicity in \eqref{eq:root} therefore gives
$\psi_d\le c$, with the asserted equality case.

At $c$, the polynomial in \eqref{eq:oldroot} has value
\[
  (c-1)f_{\bw}(c)-w_0c^{d-1}
    =w_0\left(\frac1{c-1}-c^{d-1}\right).
\]
It vanishes for constant profiles and is negative otherwise.
The uniqueness of the source's positive root
\cite[Remark~4(1)]{BKKKZ} gives the comparison with $\rho_{\bw}$.
Finally, for $d\ge2$ one has $f_{\bw}(2)>w_0$, hence $c<2$.
For $d=1$, equation~\eqref{eq:root} gives $c=2$.
\end{proof}

The rational examples in \cite[Theorem~5]{BKKKZ} are governed by the
largest positive root $r_{\bw}$ of
\[
  (x-1)f_{\bw}(x)-x^{d-1}.
\]
Because every weight is at least one, this polynomial is at least
$x^d-x^{d-1}-1$ for $x>1$. It is strictly larger there if some weight
exceeds one. Thus $r_{\bw}\le\psi_d\le c$, with $r_{\bw}<\psi_d$
if a weight exceeds one. When all weights are one, the examples have
finite $\lim a_n^{1/\psi_d^n}$ and do not meet the infinite-limsup
hypothesis here. For $(2,1)$ the familiar gap between $\sqrt2$ and
$\sqrt3$ remains. For $(3,2,1)$ the rational-example polynomial is
$x^3+x-3$, whereas the sufficient base here is $\sqrt[3]{4}$.
The present theorem does not close these gaps.

One possible way to apply an earlier criterion is to replace $a_n$
by a finite product of its shifts. The following comparison treats
those substitutions explicitly.

\begin{proposition}\label{prop:products}
Suppose
\[
  f_{\bw}(x)=V(x)G(x),\qquad
  V(x)=\sum_{i=0}^{k-1}v_ix^i,\qquad G(x)=\sum_{j=0}^g q_jx^j,
\]
where the coefficients are nonnegative integers,
$v_0q_0=w_0>0$, $v_{k-1}>0$ and $q_g>0$.
Let $v_{\max}=\max_i v_i$, and let $\rho_v$ be the sufficient base in
\cite[Theorem~3]{BKKKZ} for the profile $(v_0,\ldots,v_{k-1})$.
Then $c_{\bw}\le\rho_v$, with equality only if both $G$ and the
coefficient profile of $V$ are constant. In particular, the inequality
is strict for every nonconstant profile $\bw$ satisfying
\eqref{eq:weights}.
\end{proposition}

To interpret the factorization, set $y_n=\prod_{j=0}^g a_{n+j}^{q_j}$.
Then $f_{\bw}=VG$ says exactly that $\prod_i y_{n+i}^{v_i}=x_n$:
the profile $v$ applied to the derived sequence $y_n$ reproduces the
original denominators.

\begin{proof}
At $x=\rho_v$ one has $(x-1)V(x)=v_{\max}x^{k-1}$.
Since $V(x)\le v_{\max}T_k(x)$ for $x>1$, it follows that
$x^{k-1}\le x^k-1$, or $(x-1)x^{k-1}\ge1$.
Consequently,
\begin{equation}\label{eq:productcomparison}
  \frac{(\rho_v-1)^2 f_{\bw}(\rho_v)}{w_0}
   =\bigl((\rho_v-1)\rho_v^{k-1}\bigr)
       \frac{v_{\max}}{v_0}\frac{G(\rho_v)}{q_0}\ge1.
\end{equation}
Monotonicity in \eqref{eq:root} proves $c_{\bw}\le\rho_v$.
Equality forces $\rho_v=\psi_k$, equality in $V\le v_{\max}T_k$,
and $G(\rho_v)=q_0$. Hence all $v_i$ are equal and $G$ is constant.
Conversely those conditions give equality and a constant profile $\bw$.
\end{proof}

For every fixed $t>1$, the bounds
\[
  q_0\log a_n\le\log y_n\le G(1)\log a_{n+g}
\]
show that $\limsup(\log y_n)/t^n=\infty$ is equivalent to the same
condition for $a_n$; the fixed shift changes the normalization by only
the factor $t^g$. Thus these direct applications of the earlier theorem
do not recover the smaller sufficient base. This conclusion concerns
the stated finite product substitutions, not every possible method of
deciding irrationality for an individual series.

\section*{Verification}
An independent audit found no gap in the stated theorems and propositions.
This was not formal verification.

\begin{thebibliography}{9}
\bibitem{BKKKZ}
K.~Barreto, J.~Kang, S.-h.~Kim, V.~Kova\v{c} and S.~Zhang,
\emph{Irrationality of rapidly converging series: A problem of
Erd\H{o}s and Graham},
Bull.\ Lond.\ Math.\ Soc.\ \textbf{58} (2026), no.~8, e70457.
\href{https://doi.org/10.1112/blms.70457}{doi:10.1112/blms.70457}.
\href{https://arxiv.org/abs/2601.21442v3}{arXiv:2601.21442v3}.

\bibitem{Erdos1975}
P.~Erd\H{o}s,
\emph{Some problems and results on the irrationality of the sum of
infinite series},
J.\ Math.\ Sci.\ \textbf{10} (1975), 1--7.
\href{https://www.renyi.hu/~p_erdos/1976-44.pdf}{Full text}.

\bibitem{ErdosGraham}
P.~Erd\H{o}s and R.~L.~Graham,
\emph{Old and new problems and results in combinatorial number theory},
Monographies de L'Enseignement Math\'ematique, vol.~28,
Universit\'e de Gen\`eve, L'Enseignement Math\'ematique, Geneva, 1980.
\href{https://mathweb.ucsd.edu/~ronspubs/80_11_number_theory.pdf}{Full text}.

\bibitem{Earlier}
twinheterodox,
\emph{A $\sqrt3$ growth criterion for weighted reciprocal series},
Hexagon preprint, 13 November 2026, 6 pp.
\href{https://hexagonmath.org/2611.00428v1}{Hexagon 2611.00428v1}.
\end{thebibliography}

\end{document}
